La motivación central de este trabajo de graduación surge de una brecha estructural verificable: las startups y empresas tecnológicas que desarrollan servicios sobre Node.js rara vez cuentan con un equipo de seguridad dedicado, mientras que la superficie de ataque de sus APIs REST --inyección SQL, Cross-Site Scripting, Path Traversal/LFI e inyección de comandos del sistema operativo-- corresponde a categorías de riesgo documentadas de forma consistente en el estándar OWASP Top~10 como las más críticas para aplicaciones web modernas \parencite{owasp2021}. En un ecosistema donde una proporción mayoritaria de las aplicaciones web depende de componentes de código abierto propensos a vulnerabilidades \parencite{snyk2024}, y donde los reportes anuales de la industria documentan de forma sostenida el bajo costo de explotación de fallos de lógica de negocio frente a los altos costos de mantenimiento de controles perimetrales \parencite{verizon2024}, la ausencia de mecanismos de detección accesibles a nivel de capa de aplicación es una condición estructural del ecosistema, y no un vacío incidental, que este proyecto busca atender directamente.

\section{Importancia del problema}

El problema central que motiva esta investigación es la ausencia de mecanismos automatizados, técnicamente fundamentados y económicamente viables para detectar ataques ampliamente documentados y estudiados en organizaciones sin capacidad de inversión en seguridad dedicada. Las soluciones existentes revisadas en el capítulo de Antecedentes --cortafuegos de aplicación web (WAF), sistemas de detección de intrusiones (IDS) y plataformas de gestión de eventos de seguridad (SIEM)-- comparten una limitación común: exigen mantenimiento especializado continuo, infraestructura de soporte o un costo de licenciamiento que las sitúa fuera del alcance de organizaciones pequeñas \parencite{verizon2024}. Esta condición no es exclusiva de un contexto geográfico particular, pero se agudiza en economías donde el crecimiento del ecosistema tecnológico no ha sido acompañado de una madurez institucional equivalente en ciberseguridad, como es el caso de Guatemala, donde los ciberataques crecieron un 200\% entre 2023 y 2024 y el país se ubica entre los diez con mayor exposición cibernética de Latinoamérica \parencite{mastercard2024}. Abordar este problema es importante porque la ausencia de protección a nivel de aplicación expone directamente datos de clientes, información financiera y la continuidad operativa de organizaciones que, por su tamaño, no pueden absorber el costo de un incidente de seguridad significativo.

\section{Relevancia técnica y científica}

Desde la perspectiva técnica, este trabajo es relevante porque atiende un vacío identificado de forma explícita en el capítulo de Antecedentes: la literatura de detección de intrusiones mediante aprendizaje automático se ha desarrollado predominantemente sobre datos de flujo de red, no sobre la semántica de peticiones HTTP individuales, y ni la revisión de arquitecturas híbridas de aprendizaje supervisado y no supervisado de \textcite{zhen2025} ni la revisión sistemática de cortafuegos de aplicación web potenciados por inteligencia artificial de \textcite{rahman2026} identifican, entre la literatura que cubren, un sistema implementado bajo una arquitectura embebida dentro del proceso de la aplicación protegida. Investigar la viabilidad de un modelo híbrido de aprendizaje automático ejecutado en tiempo de inferencia dentro de un servidor Node.js, sujeto a restricciones estrictas de latencia impuestas por su modelo de concurrencia de un solo hilo, constituye un aporte técnico concreto: extiende una disciplina de diseño --la inferencia embebida bajo restricciones de recursos, consolidada en la literatura de Edge AI para hardware de microcontroladores-- hacia un dominio de ejecución distinto para el cual no se identificó antecedente documentado. Este componente de novedad metodológica es lo que justifica, desde el punto de vista científico, el carácter aplicado-experimental de la investigación: su meta es generar evidencia empírica sobre si dicha combinación arquitectónica es técnicamente viable, en lugar de implementar una solución cuya viabilidad ya esté validada de antemano.

\section{Relevancia social, económica e institucional}

En el plano social y económico, la relevancia del proyecto se sostiene en la magnitud documentada del problema en el contexto local. El ecosistema de startups guatemaltecas exhibe una de las tasas de actividad emprendedora más altas del mundo, pero opera en un entorno donde la ausencia de un marco legal específico de ciberseguridad y la escasez de profesionales especializados en la región limitan estructuralmente la capacidad de estas organizaciones para invertir en controles de seguridad tradicionales. Una herramienta distribuida sin costo a través de un registro de paquetes de uso extendido, que no requiere modificaciones a la infraestructura existente ni conocimiento especializado para su integración, reduce directamente las dos barreras de adopción más citadas en la literatura de la industria: el presupuesto disponible y la complejidad técnica de implementación \parencite{verizon2024}. En el plano institucional, el proyecto es relevante para la Universidad del Valle de Guatemala en la medida en que constituye una contribución documentada y reproducible en la intersección entre ciencia de datos, ciberseguridad e ingeniería de software, un área de creciente relevancia curricular y de investigación aplicada dentro de la Facultad de Ingeniería.

\section{Beneficiarios de la investigación}

Los beneficiarios directos de este trabajo son, en primer lugar, las startups y empresas tecnológicas que operan servicios sobre Node.js sin un equipo de seguridad dedicado, para quienes la investigación busca proveer una herramienta de protección accesible a nivel de capa de aplicación. En segundo lugar, son beneficiarios directos los propios autores del proyecto integrador, para quienes el trabajo representa la aplicación rigurosa de las competencias desarrolladas en sus respectivas áreas de formación --inteligencia artificial y aprendizaje automático, en el caso del componente documentado en este protocolo, y ciberseguridad, en el caso del componente desarrollado en paralelo-- sobre un problema real y técnicamente exigente. Como beneficiarios indirectos se identifican, por un lado, la comunidad académica y estudiantil de la Universidad del Valle de Guatemala, que podrá disponer de la metodología documentada de construcción de un pipeline de detección de amenazas sobre arquitecturas orientadas a eventos como referencia para cursos o trabajos de graduación posteriores; y por otro, de manera más amplia, los usuarios finales de las aplicaciones que eventualmente adopten la herramienta, cuya información personal y financiera se beneficia de manera transitiva de cualquier mejora en la postura de seguridad de los servicios que utilizan.

\section{Síntesis}

La justificación de este trabajo de graduación se apoya en la combinación específica de tres factores que, tomados en conjunto, delimitan un problema real y no resuelto, y no en la novedad de los ataques que busca detectar, ya ampliamente catalogados en estándares como OWASP Top~10 \parencite{owasp2021}: una brecha de adopción de controles de seguridad documentada de forma consistente en la literatura de la industria \parencite{verizon2024, snyk2024}; un vacío metodológico identificado en la literatura académica respecto a la detección de amenazas embebida en tiempo de ejecución dentro de aplicaciones Node.js; y un contexto local, el ecosistema de startups guatemaltecas, donde ambos factores convergen con particular intensidad \parencite{mastercard2024}. Este capítulo no pretende anticipar en qué medida la solución propuesta logrará resolver dicho problema; ese es precisamente el objeto de la investigación desarrollada en los capítulos siguientes. Lo que este capítulo establece es que el problema es real, que su abordaje es pertinente tanto desde una perspectiva técnica como social e institucional, y que existen beneficiarios concretos, directos e indirectos, para quienes una respuesta fundamentada a este problema representa un aporte de valor verificable.
